\section{Conclusiones y Líneas Abiertas}

El sistema desarrollado demuestra la viabilidad técnica de clasificar perfiles de riesgo mediante redes neuronales convolucionales con una alta tasa de éxito. La integración de \textbf{MTCNN} y \textbf{Grad-CAM} ha sido crucial para garantizar que el modelo aprenda rasgos faciales genuinos y no artefactos del fondo de las imágenes.

Como conclusión principal, la arquitectura \textbf{ResNet50} ha mostrado un equilibrio superior entre precisión y estabilidad, mientras que \textbf{VGG-16} presenta mayor sensibilidad a detalles específicos. El balanceo de datos mediante \textbf{DeepFace} ha sido la pieza angular para garantizar un sistema éticamente responsable y libre de sesgos raciales directos.

\subsection{Líneas Futuras de Investigación}
\begin{itemize}
    \item \textbf{Análisis Multimodal:} Integrar registros de audio o texto forense para complementar el análisis visual.
    \item \textbf{Sistemas de Tiempo Real:} Optimizar el pipeline para su ejecución en dispositivos móviles con \textit{MobileNetV3}.
    \item \textbf{Debate Ético y Legal:} Profundizar en las implicaciones del uso de estas herramientas en sistemas judiciales reales y su cumplimiento con el RGPD.
\end{itemize}

\section{Referencias}

\begin{enumerate}
    \item Álvarez Navarro, A., et al. (2022). \textit{Revisión de los métodos de reconocimiento facial en imágenes RGB-D}. Revista Cubana de Ciencias Informáticas.
    \item Fructuoso Freire Montero, A. (2022). \textit{El reconocimiento facial como instrumento de investigación y prevención del delito}. Anuario da Facultade de Dereito da Universidade da Coruña.
    \item Garcia Silva, C. E., \& Mazon Loaiza, J. D. (2024). \textit{El reconocimiento facial como instrumento de investigación y prevención del delito}. Anatomía Digital.
    \item Goring, C. (1913). \textit{The English Convict: A Statistical Study}. His Majesty's Stationery Office.
    \item He, K., Zhang, X., Ren, S., \& Sun, J. (2016). \textit{Deep Residual Learning for Image Recognition}. CVPR.
    \item Kranthi Kumar, K., et al. (2021). \textit{Criminal face identification system using deep learning algorithm multi-task cascade neural network (MTCNN)}. Materials Today: Proceedings.
    \item Lombroso, C. (1876). \textit{L'Uomo Delinquente}. Hoepli.
    \item Rani, R., et al. (2025). \textit{Face Recognition System for Criminal Identification in CCTV Footage Using Keras and OpenCV}. Ingénierie des Systèmes d'Information.
    \item Selvaraju, R. R., et al. (2017). \textit{Grad-CAM: Visual Explanations from Deep Networks via Gradient-based Localization}. ICCV.
    \item Simonyan, K., \& Zisserman, A. (2014). \textit{Very Deep Convolutional Networks for Large-Scale Image Recognition}.
    \item Zhang, K., Zhang, Z., Li, Z., \& Qiao, Y. (2016). \textit{Joint Face Detection and Alignment using Multi-task Cascaded Convolutional Networks}. IEEE Signal Processing Letters.
\end{enumerate}
\clearpage 